\section*{Chapter \ref{ch:g9:price-controls-taxes-subsidies} --- Price Controls, Taxes and Subsidies}

\begin{solution}{exo:g9:price-controls-taxes-subsidies:1}
They reduced it by $15\,\%$. The authors calculate that rents across the city were about $5.1\,\%$
higher; this is a calculation from their estimates (the supply cut and an assumed response of renters),
not a measured effect.
\end{solution}

\begin{solution}{exo:g9:price-controls-taxes-subsidies:2}
Demanded $120 - 36 = 84$, supplied $20 + 54 = 74$: a \omterm{def:g7:markets-and-prices:shortage}{shortage} of 10 units.
\end{solution}

\begin{solution}{exo:g9:price-controls-taxes-subsidies:3}
Demanded $120 - 44 = 76$, supplied $20 + 66 = 86$: 10 units unsold.
\end{solution}

\begin{solution}{exo:g9:price-controls-taxes-subsidies:4}
$13 - 5 = 8$ dalens.
\end{solution}

\begin{solution}{exo:g9:price-controls-taxes-subsidies:5}
$Q^s = 20 + 6(p - 10) = -40 + 6p$. Then $120 - 4p = -40 + 6p$, so $160 = 10p$: buyers pay 16 dalens,
sellers keep $16 - 10 = 6$ dalens, and $120 - 64 = 56$ units are sold.
\end{solution}

\begin{solution}{exo:g9:price-controls-taxes-subsidies:6}
$10 \times 56 = 560$ dalens.
\end{solution}

\begin{solution}{exo:g9:price-controls-taxes-subsidies:7}
Nothing: the ceiling is above the \omterm{def:g7:demand-and-supply-first-look:equilibrium-price}{equilibrium price} of 10 dalens, so the \omterm{def:g7:money:price}{price} stays at 10 and 80 units are
sold.
\end{solution}

\begin{solution}{exo:g9:price-controls-taxes-subsidies:8}
$Q^s = 20 + 6(p + 2) = 32 + 6p$. Then $120 - 4p = 32 + 6p$, so $88 = 10p$: buyers pay $8.80$ dalens,
sellers receive $10.80$, and $120 - 35.2 = 84.8$ units are sold. Cost: $2 \times 84.8 = 169.60$ dalens.
\end{solution}

\begin{solution}{exo:g9:price-controls-taxes-subsidies:9}
For example: the maximum rent is above the rent the \omterm{def:g7:markets-and-prices:market}{market} would set anyway, so it changes nothing; or
tenants pay extra in other ways that the law does not see (a fee, a payment for furniture), so the real
rent is higher than the legal one.
\end{solution}

\begin{solution}{exo:g9:price-controls-taxes-subsidies:10}
Without the \omterm{def:g8:taxes-and-public-spending:tax}{tax}: $100 - p = 4p$, so $p = 20$ and $Q = 80$. With the \omterm{def:g8:taxes-and-public-spending:tax}{tax}: $Q^s = 4(p - 5) = 4p - 20$, so
$100 - p = 4p - 20$, $120 = 5p$, $p = 24$: buyers pay 24 dalens, sellers keep 19, and 76 units are sold.
Buyers bear 4 of the 5 dalens: they respond less to the \omterm{def:g7:money:price}{price} ($b = 1$) than sellers ($d = 4$).
\end{solution}

\begin{solution}{exo:g9:price-controls-taxes-subsidies:11}
With $Q^s = 20 + 6p$: the \omterm{def:g7:money:price}{price} rises by $5 \times 6 \div 10 = 3$ dalens. With $Q^s = 70 + p$: before,
$120 - 4p = 70 + p$ gives $p = 10$; with the \omterm{def:g8:taxes-and-public-spending:tax}{tax}, $120 - 4p = 70 + (p - 5)$ gives $55 = 5p$, $p = 11$: the
buyers' \omterm{def:g7:money:price}{price} rises by only 1 dalen ($5 \times 1 \div 5$). Sellers who respond little bear most of the
\omterm{def:g8:taxes-and-public-spending:tax}{tax}.
\end{solution}

\begin{solution}{exo:g9:price-controls-taxes-subsidies:12}
Many other things changed in the city after 1994: \omterm{def:g7:households:income}{incomes}, the number of people who wanted to live there,
rents everywhere. A before-and-after comparison would mix all of them with the effect of rent control.
Similar buildings that were not covered went through the same changes, so the difference between the two
groups is due to rent control.
\end{solution}

\begin{solution}{pb:g9:price-controls-taxes-subsidies:1}
\textbf{1.} At 2 dalens: $7\,000 - 2\,000 = 5\,000$ loaves; at 3: $4\,000$.
\textbf{2.} At 2 dalens: $-2\,000 + 4\,000 = 2\,000$ loaves; at 3: $4\,000$.
\textbf{3.} $7\,000 - 1\,000p = -2\,000 + 2\,000p$, so $9\,000 = 3\,000p$: 3 dalens and $4\,000$ loaves.
\textbf{4.} $3 \times 4\,000 = 12\,000$ dalens.
\textbf{5.} Demanded $7\,000 - 2\,500 = 4\,500$; supplied $-2\,000 + 5\,000 = 3\,000$; \omterm{def:g7:markets-and-prices:shortage}{shortage} $1\,500$
loaves.
\textbf{6.} $3\,000$ loaves, the quantity bakers offer.
\textbf{7.} $7\,000 - 1\,000p = 3\,000$ gives $p = 4$ dalens.
\textbf{8.} For example: queues before dawn; smaller or poorer loaves; bread kept for regular customers; some
bakers closing for the winter.
\textbf{9.} $4\,500$ loaves.
\textbf{10.} $-2\,000 + 2\,000p = 4\,500$ gives $p = 3.25$ dalens.
\textbf{11.} $3.25 - 2.50 = 0.75$ dalen per loaf.
\textbf{12.} $0.75 \times 4\,500 = 3\,375$ dalens a day.
\textbf{13.} Families pay $2.50$ instead of 3 dalens; bakers receive $3.25$ instead of 3; $4\,500$ loaves
are baked instead of $4\,000$. The town's taxpayers pay the difference, $3\,375$ dalens a day.
\end{solution}
